\subsection{Tree-Quality Audit}
\label{sec:tree_quality_audit}
We next examine whether node difficulty is associated with measurable tree properties. For each node, we combine branching factor, split entropy, largest-child share, and depth with the traversal accuracy and model-specific grounding measurements reported above. Nodes are identified by their root-to-node answer paths. For an internal node $v$ with children $j=1,\dots,d_v$, let $C_{vj}$ be the classes reachable under child $j$ and $p_{vj}=|C_{vj}|/\sum_k |C_{vk}|$; split entropy is $H_v=-\sum_j p_{vj}\log_2 p_{vj}$, reported in bits and \emph{not} normalized by $\log_2 d_v$, since branching factor enters the model as a separate predictor.

The four trees contain 128 internal nodes and 153 leaves (GTSRB 22/43, MNIST 8/10, CIFAR-10 9/10, MSCOCO 89/90). Every internal node has outdegree at least two --- GTSRB $d{=}2$ (16 nodes), $3$ (2), $5$ (2), $8$ (2); MNIST $d{=}2$ (7), $3$ (1); CIFAR-10 and MSCOCO binary throughout --- so $L=1+\sum_v (d_v-1)$ holds exactly in each tree and no unary nodes are present. MSCOCO has 90 leaves because our label space spans the full 1--90 COCO identifier range rather than the 80 annotated categories (Section~\ref{sec:datasets}); the same 90-identifier set defines the zero-shot candidate list, so both methods choose from an identical option set. Five MSCOCO nodes at depths 4--6 are never reached \emph{along a fully correct traversal path} and carry no measured accuracy; the models do visit them, but only after an earlier error. One separates two unannotated identifiers and is unmeasurable by construction; the other four sit above \emph{potted plant} and \emph{dining table}, whose traversals diverge at depths 0--2 in every instance. This leaves 123 nodes with measured accuracy and grounding gaps. Split entropy ($r=-0.24$) and branching factor ($r=-0.17$) correlate negatively with traversal accuracy, whereas depth has a weak positive correlation ($r=0.18$); the depth association reverses within three of the four trees (CIFAR-10 $-0.40$, MNIST $-0.12$, MSCOCO $-0.08$, GTSRB $+0.56$) and so reflects between-dataset differences. Evenly split nodes (largest child $\leq60\%$ of reachable classes, $n=68$) average 81.6\% accuracy, compared with 88.2\% for more lopsided nodes ($n=55$).

To assess grounding without inflating its association with dataset difficulty, we pair each model's grounding gap with its own traversal accuracy. Of the $123\times3=369$ possible pairings, 347 are observed: GPT-4o reached all 123 nodes, Qwen-VL-Max 122 and LLaMA-3.2 only 103, and one pairing lacks a grounding measurement. This missingness is informative --- the 20 nodes LLaMA never reached average depth 8.7 against 4.6 for those it did --- so the weakest model is under-represented where trees are deepest. The outcome is node accuracy, a proportion in $[0,1]$; all four predictors are standardized to zero mean and unit variance, and coefficients are unweighted. The regressions include model fixed effects; LLaMA scores 15.5 points below GPT-4o ($95\%$ CI $[-21.1,-9.9]$, $p<0.0001$). We report the naive fit ($R^2=0.16$) alongside specifications with errors clustered by node and a node-level random intercept. We do not report errors clustered by dataset: with only four clusters such inference is unreliable, and in our data it is anti-conservative, shrinking the grounding-gap standard error below its naive value.

Split entropy is the most robust predictor, remaining negative and significant in all three specifications ($\beta=-0.09$, $95\%$ CI $[-0.13,-0.04]$, $p\leq0.007$). Depth is significant in all three ($\beta=+0.035$, $p=0.004$--$0.039$), branching factor in two ($\beta=+0.052$, $p=0.029$--$0.11$; note that its sign reverses relative to the bivariate correlation once entropy is controlled, reflecting their collinearity). The grounding gap is consistently negative ($\beta=-0.03$) and significant in two specifications, but not in the node-random-intercept model ($p=0.23$), which relies mainly on within-node variation across models. Because node accuracies rest on very unequal denominators (1 to 1000 decisions, median 18.5), we refit the naive specification weighted by each node's decision count: split entropy ($p=0.0001$) and the grounding gap ($p=0.0001$) remain significant, whereas depth ($p=0.37$) and branching factor ($p=0.07$) do not.

Thus, nodes with larger grounding gaps are harder on average, but within a node, the model with the larger gap is not reliably more error-prone. A pooled analysis would overstate the grounding association ($\beta=-0.07$, $p<0.0001$), whereas the negative relationship between split entropy and accuracy is the one structural effect that survives controlling for the other factors, node-level dependence, and decision-count weighting. Depth and branching factor are less stable: both lose significance under weighting, and the apparent depth effect is partly attributable to between-tree differences and to informative missingness. We therefore read tree structure as a weak but consistent correlate of node difficulty rather than a reliable predictor, and we limit these conclusions to the four trees evaluated here.
